\documentclass[aspectratio=169,11pt]{beamer}

% White, minimal slides
\usetheme{default}
\usecolortheme{default}
\setbeamertemplate{navigation symbols}{}
\setbeamertemplate{footline}[frame number]
\setbeamercolor{background canvas}{bg=white}
\setbeamercolor{normal text}{fg=black,bg=white}
\setbeamercolor{frametitle}{fg=black,bg=white}
\setbeamercolor{title}{fg=black,bg=white}
\setbeamercolor{block title}{fg=black,bg=white}
\setbeamercolor{block body}{fg=black,bg=white}
\setbeamertemplate{blocks}[default]
\setbeamertemplate{itemize item}{--}
\setbeamertemplate{itemize subitem}{--}

\usepackage{amsmath, amssymb, bm}
\usepackage{booktabs}
\usepackage{array}
\usepackage{mathtools}
\usepackage{hyperref}

\newcommand{\R}{\mathbb{R}}
\newcommand{\KL}{D_{\mathrm{KL}}}
\newcommand{\diag}{\operatorname{diag}}
\newcommand{\Had}{\mathrm{Had}}

\title{From Learned Rotations to Search-Based Hadamard Rotations}

\author{}
\date{}

\begin{document}

% 1
\begin{frame}
  \titlepage
\end{frame}

% 2
\begin{frame}{why SpinQuant works}
\vspace{0.2cm}
LLM quantization is hard because a few large coordinates dominate the quantization range.

\vspace{0.3cm}
For a tensor $X$, uniform quantization roughly uses
\[
\alpha \propto \frac{\max |X|}{2^{b-1}-1}.
\]
A large outlier increases $\alpha$, so most ordinary values are represented with fewer effective levels.

\vspace{0.4cm}
\textbf{Rotation idea:}
\[
XW = (XR)(R^{-1}W), \qquad R^\top R = I.
\]
In full precision, the model is unchanged. After quantization, the rotated tensors may have fewer outliers and smaller quantization error.
\end{frame}

% 3
\begin{frame}{Small example: rotation spreads an outlier}
Consider four activation channels with a large value only in the first coordinate:
\[
 x = \begin{bmatrix} 10 & 0 & 0 & 0 \end{bmatrix}.
\]
Use the normalized $4\times 4$ Hadamard rotation
\[
R = \frac{1}{2}
\begin{bmatrix}
1&1&1&1\\
1&-1&1&-1\\
1&1&-1&-1\\
1&-1&-1&1
\end{bmatrix}.
\]
Then
\[
 xR = \begin{bmatrix}5&5&5&5\end{bmatrix}.
\]

\vspace{0.2cm}
The norm is unchanged:
\[
\|x\|_2 = 10, \qquad \|xR\|_2 = 10,
\]
but the maximum coordinate drops:
\[
\|x\|_\infty = 10, \qquad \|xR\|_\infty = 5.
\]
\end{frame}

% 4
% \begin{frame}{Why random rotations already help}
% For a high-dimensional vector $x\in\R^d$, a random orthogonal rotation tends to spread energy across coordinates.

% \vspace{0.3cm}
% If
% \[
%  x = (10,0,\ldots,0),
% \]
% then after a random orthogonal rotation $R$,
% \[
% \|xR\|_2 = 10,
% \]
% but a typical coordinate has scale
% \[
% \frac{\|x\|_2}{\sqrt d} = \frac{10}{\sqrt d}.
% \]

% \vspace{0.3cm}
% So random rotations often reduce coordinate outliers. The limitation is variance: different random rotations can give noticeably different quantized accuracy. SpinQuant uses this observation to motivate learning the rotation.
% \end{frame}

% 5
\begin{frame}{What SpinQuant optimizes}
SpinQuant inserts rotations into the Transformer in places that preserve the full-precision function.

\vspace{0.2cm}
Two important learned rotations:
\begin{itemize}
  \item $R_1$: residual-stream rotation, mergeable into nearby weights.
  \item $R_2$: attention value/output rotation, also mergeable.
\end{itemize}

\vspace{0.2cm}
The optimization problem is
\[
\min_{R_1,R_2\in\mathcal M} \mathcal L_Q(R_1,R_2\mid W,X),
\]
where
\[
\mathcal M = \{R: R^\top R = I\}.
\]

\vspace{0.2cm}
SpinQuant uses Cayley SGD to keep the rotations orthonormal while minimizing the quantized-model loss.
\end{frame}

% 6
\begin{frame}{Motivation: can we avoid dense Cayley optimization?}
SpinQuant is accurate, but the optimization can be costly.

and very sensitive to the learning rate.

\vspace{0.2cm}
Potential concerns:
\begin{itemize}
  \item We need large GPU memory before to produce the quantised model.
  \item Very sensitive to the learning rates.
  % \item Backpropagation through quantized forward passes needs calibration activations and gradients.
  % \item Cayley-style updates preserve orthogonality but are more complex than simple sign or permutation search.
\end{itemize}

\vspace{0.3cm}
% \textbf{Research question:}

Can we get much of the benefit of learned rotations using a cheaper, discrete, search-based fast-Hadamard family?
\end{frame}

% 7
\begin{frame}{EvoPress: the search-based pruning inspiration}
EvoPress studies dynamic LLM compression: choose different compression levels for different layers while satisfying a global budget.

% \vspace{0.2cm}
% Instead of assuming additive layer importance,
% \[
% \text{total error} \approx \sum_{\ell} \text{layer error}_{\ell},
% \]
EvoPress evaluates whole compressed candidates.

\vspace{0.2cm}
formulation:
\[
\hat M^* = \arg\min_{v\in[m]^n} \KL(P_M\,\\|\,Q_{\hat M_v})
\quad\text{s.t.}\quad g(\hat M_v)\le C.
\]

\vspace{0.2cm}
Key algorithmic idea:
\begin{itemize}
   \item search space is large
  \item Maintain one parent compression profile.
  \item Generate local budget-preserving mutations.
  \item Use multi-step selection: many cheap evaluations first, few expensive evaluations later.
\end{itemize}
\end{frame}

% 8
% \begin{frame}{Key lesson from EvoPress}
% EvoPress is useful because the search space is huge, but candidate evaluation can be staged.

% \vspace{0.3cm}
% \begin{center}
% \begin{tabular}{>{\raggedright\arraybackslash}p{0.33\linewidth} >{\raggedright\arraybackslash}p{0.52\linewidth}}
% \toprule
% Problem & EvoPress design choice \\
% \midrule
% Huge discrete space & Local mutations, not exhaustive search \\
% Expensive model evaluation & Multi-step selection with increasing tokens \\
% Layer interactions & Evaluate whole candidate models \\
% Global compression budget & Mutations preserve the budget \\
% \bottomrule
% \end{tabular}
% \end{center}

% \vspace{0.3cm}
% This suggests a similar strategy for rotation search: do not optimize dense matrices; search over structured rotations and evaluate candidates cheaply before full quantization.
% \end{frame}

% 9
\begin{frame}{search over structured Hadamard rotations?}

  \vspace{0.2cm}
A valid Hadamard matrix must satisfy
\[
HH^\top = dI.
\]
Most random sign matrices are not Hadamard.

\vspace{0.2cm}
Instead, search over fast structured families such as
\[
R(s) = \frac{1}{\sqrt d} D(s)H,
\qquad
D(s)=\diag(s_1,\ldots,s_d),\quad s_i\in\{\pm1\}.
\]
Now the search space is
\[
2^d
\]
not
\[
2^{d^2}.
\]

\vspace{0.2cm}
Richer but still fast families are possible:
\[
R = D_0 H D_1 H D_2 \cdots H D_k.
\]
\end{frame}

% 10
\begin{frame}{Why sign-Hadamard search may help}
For a signed Hadamard rotation,
\[
(xR(s))_j = \frac{1}{\sqrt d}\sum_i x_i s_i H_{ij}.
\]
The signs $s_i$ determine which channels cancel or reinforce after mixing.

\vspace{0.2cm}
Random signs already spread energy, but a searched sign pattern may better reduce:
\[
\max_j |(xR(s))_j|,
\]
or the actual fake-quantization loss:
\[
\|Q(xR(s)) - xR(s)\|^2.
\]

\vspace{0.2cm}
The idea is strongest when LLM outliers are structured across channels, rather than isolated single-coordinate spikes.
\end{frame}

% 11
\begin{frame}{Search over layers}
\textbf{Candidate:} sign vectors, permutations, or multiple signed-Hadamard factors for each rotation location.
\[
\theta = \{s_{R_1}, s_{R_2}^{(1)},\ldots,s_{R_2}^{(L)}, s_{R_3}, s_{R_4}\}.
\]

\vspace{0.2cm}
\textbf{Search objective:}
\[
\min_{\theta}\; \KL(P_{\mathrm{FP}}\,\|\,P_{\mathrm{fake\text{-}quant}}(\theta)).
\]

\vspace{0.2cm}
% \textbf{Multi-stage evaluation:}
% \begin{enumerate}
%   \item Cheap proxy: activation outlier score or quantization error on cached activations.
%   \item Medium proxy: fake-quantized KL on a small token batch.
%   \item Final selection: larger calibration set; then run final quantization once.
% \end{enumerate}
\end{frame}

% 12
\begin{frame}{Methodology: first implementation target}
\textbf{Goal:} start with the smallest search problem that still tests the core idea.

\vspace{0.2cm}
\textbf{Search sites:}
\begin{itemize}
  \item Search only the per-layer attention value/output rotations $R_2^{(\ell)}$.
  \item Keep the global residual rotation $R_1$ fixed at a random Hadamard-style initialization.
  \item This isolates whether search can improve over random Hadamard without the harder global search over $R_1$.
\end{itemize}

\vspace{0.2cm}
\textbf{Rotation family per site:}
\[
R_2^{(\ell)} = D_0^{(\ell)} H D_1^{(\ell)},
\qquad
D_0^{(\ell)}, D_1^{(\ell)} \in \{\pm1\}^{d_{\mathrm{head}}\times d_{\mathrm{head}}}
\]
where each $D$ is diagonal with $\pm 1$ entries.

\vspace{0.2cm}
This uses one fast Hadamard mixing step per site, but still gives a searchable family through the sign patterns.
\end{frame}

% 13
\begin{frame}{Methodology: search operators}
\textbf{Per-site evolutionary search:}
\begin{itemize}
  \item Maintain a small population of candidate sign patterns for one site at a time.
  \item Evaluate, select top candidates, mutate them, and repeat for a fixed number of generations.
  \item Move layer by layer: optimize one $R_2^{(\ell)}$ while holding all other sites fixed.
\end{itemize}

\vspace{0.2cm}
\textbf{Candidate encoding:}
\[
\phi^{(\ell)} = \bigl(s_0^{(\ell)}, s_1^{(\ell)}\bigr),
\qquad
s_0^{(\ell)}, s_1^{(\ell)} \in \{\pm1\}^{d_{\mathrm{head}}}
\]
which defines the diagonals of $D_0^{(\ell)}$ and $D_1^{(\ell)}$.

\vspace{0.2cm}
\textbf{Mutation examples:}
\begin{itemize}
  \item Flip a small random subset of signs in $s_0^{(\ell)}$ or $s_1^{(\ell)}$.
  \item Apply multiple mutation strengths, e.g. small flips for local refinement and larger flips for exploration.
  \item Optionally keep one elite candidate unchanged each generation.
\end{itemize}
\end{frame}

% 14
\begin{frame}{Methodology: initialization}
\textbf{Search one site at a time.}

\vspace{0.2cm}
For layer $\ell$, keep all other rotations fixed and optimize
\[
R_2^{(\ell)} = D_0^{(\ell)} H D_1^{(\ell)}.
\]

\vspace{0.2cm}
\textbf{Initial population for site $\ell$:}
\begin{itemize}
  \item Include the current deployed candidate as the first parent.
  \item If this is the first sweep, use a random signed-Hadamard initialization.
  \item Add $P-1$ additional candidates by randomizing the sign vectors
  \[
  s_0^{(\ell)}, s_1^{(\ell)} \in \{\pm 1\}^{d_{\mathrm{head}}}.
  \]
  \item Optionally bias a few initial candidates to be small perturbations of the current parent rather than fully random.
\end{itemize}

\vspace{0.2cm}
\textbf{Why this initialization?}
\begin{itemize}
  \item It preserves a strong random-Hadamard baseline.
  \item It gives immediate diversity without requiring gradients.
  \item It supports repeated layer sweeps: later passes can start from the best candidate found so far.
\end{itemize}
\end{frame}

% 15
\begin{frame}{Methodology: mutation operators}
\textbf{Goal:} combine local refinement and occasional exploration.

\vspace{0.2cm}
For a parent candidate $\phi^{(\ell)} = (s_0^{(\ell)}, s_1^{(\ell)})$, generate offspring by:
\begin{itemize}
  \item choosing one of the two sign vectors $s_0^{(\ell)}$ or $s_1^{(\ell)}$,
  \item flipping a random subset of entries,
  \item and optionally mutating both vectors in the same offspring with small probability.
\end{itemize}

\vspace{0.2cm}
\textbf{First version: mutation-only evolution.}
\begin{itemize}
  \item No crossover between parent candidates.
  \item Each offspring is produced by mutating a single parent.
  \item This keeps the algorithm simple and makes the effect of each mutation easier to analyze.
  \item Most offspring mutate only one of $s_0^{(\ell)}$ or $s_1^{(\ell)}$; a small probability is reserved for mutating both in the same offspring.
\end{itemize}

\vspace{0.2cm}
\textbf{Mutation strengths:}
\begin{itemize}
  \item Small mutation: flip $1$ to $2$ signs.
  \item Medium mutation: flip $4$ to $8$ signs.
  \item Large mutation: flip a larger fraction of the vector for exploration.
\end{itemize}

\vspace{0.2cm}
\textbf{Practical schedule:}
\begin{itemize}
  \item Most offspring use small or medium mutations.
  \item A minority use large mutations to escape poor local optima.
  \item Keep mutation sizes fixed initially; later one can adapt them based on search progress.
\end{itemize}
\end{frame}

% 16
\begin{frame}{Methodology: selection and update rule}
\textbf{Per generation for site $\ell$:}
\begin{enumerate}
  \item Evaluate every candidate in the site population with the local fake-quant fitness.
  \item Rank candidates by lower block-output reconstruction error.
  \item Keep the top $E$ elites unchanged.
  \item Sample parents from the top fraction of the population and generate offspring until the population is refilled.
  \item Replace the site rotation with the best candidate of the new generation.
\end{enumerate}

\vspace{0.2cm}
\textbf{Parent sampling rule:}
\begin{itemize}
  \item Keep the population size fixed at $P$.
  \item Keep $E=1$ or $2$ elites unchanged.
  \item Define a parent pool as the top $25\%-50\%$ of candidates in the current generation.
  \item Sample parents from this pool with probability biased toward better fitness.
\end{itemize}

\vspace{0.2cm}
\textbf{Stopping rule for a site:}
\begin{itemize}
  \item stop after a fixed number of generations, or
  \item stop early if the best fitness has not improved for $T$ consecutive generations.
\end{itemize}

\vspace{0.2cm}
\textbf{Layer schedule:}
\begin{itemize}
  \item sweep through layers from shallow to deep,
  \item then optionally do a second full sweep using the updated site rotations as the new parents.
\end{itemize}
\end{frame}

% 17
\begin{frame}{Methodology: fitness for the first version}
\textbf{Primary local fitness: normalized block-output reconstruction error under fake quantization.}

\vspace{0.2cm}
For a candidate $R_2^{(\ell)}$, run the block with the rotated value/output pathway and fake quantization, then compare the block output to the floating-point block output:
\[
\mathcal F_{\ell}(R_2^{(\ell)}) =
\frac{
\left\|
Y_{\ell}^{\mathrm{FP}} -
Y_{\ell}^{\mathrm{fake\text{-}quant}}(R_2^{(\ell)})
\right\|_2^2
}{
\left\|
Y_{\ell}^{\mathrm{FP}}
\right\|_2^2 + \varepsilon
}.
\]

\vspace{0.2cm}
\textbf{Why use this first?}
\begin{itemize}
  \item It is much cheaper than end-to-end ASR evaluation.
  \item It is more faithful than simple outlier or activation-only proxies.
  \item It directly tests whether the searched rotation reduces quantization damage at the block level.
  \item The normalization makes scores more comparable across layers with different activation scales.
\end{itemize}
\end{frame}

% 18
\begin{frame}{Methodology: recommended first hyperparameters}
\textbf{Start simple and stable.}

\vspace{0.2cm}
Suggested initial choices:
\begin{itemize}
  \item Population size $P$: small, e.g. $8$ to $16$ candidates per site.
  \item Elites $E$: keep the top $1$ to $2$ candidates.
  \item Parent pool: top $25\%-50\%$ of the population.
  \item Generations per site: modest, e.g. $10$ to $30$.
  \item Site sweeps: start with one forward sweep, then add a second sweep if helpful.
  \item Calibration minibatch for local scoring: small fixed subset for reproducibility and low cost.
\end{itemize}

\vspace{0.2cm}
\textbf{Reasoning:}
\begin{itemize}
  \item We want to validate the search mechanism before scaling up budget.
  \item A small population keeps evaluation cost under control.
  \item Repeated sweeps can recover some of the benefit of larger populations.
\end{itemize}
\end{frame}

% 19
\begin{frame}{Methodology: staged evaluation and later extension}
\textbf{Phase 1: stable local search}
\begin{itemize}
  \item Optimize each $R_2^{(\ell)}$ using block-output reconstruction error only.
  \item Sweep across layers for one or more passes.
  \item Measure whether this beats random Hadamard and best-of-$N$ random baselines.
\end{itemize}

\vspace{0.2cm}
\textbf{Phase 2: add periodic global reranking}
\begin{itemize}
  \item Every few generations or after each layer sweep, collect the top few full-model candidates.
  \item Re-rank them using end-to-end fake-quantized CTC loss on calibration audio:
  \[
  \mathcal L_{\mathrm{CTC}}^{\mathrm{fake\text{-}quant}}(\theta).
  \]
  \item This lets cheap local fitness drive most of the search while the task-level objective corrects for cross-layer interactions.
\end{itemize}

\vspace{0.2cm}
\textbf{Longer-term extension:}
\begin{itemize}
  \item After validating $R_2$ search, extend the same framework to the harder global rotation $R_1$ using a richer family such as
  \[
  R_1 = D_0 H D_1 H D_2.
  \]
\end{itemize}
\end{frame}

% 20
\begin{frame}{Methodology: reusable search module}
\textbf{Design goal:} the search algorithm should work for both Whisper and Parakeet.

\vspace{0.2cm}
\textbf{Split the system into:}
\begin{itemize}
  \item \textbf{Search engine}: population, elites, parent sampling, mutation, stopping rule.
  \item \textbf{Model adapter}: model-specific site discovery and block replay logic.
  \item \textbf{Site evaluator}: computes the local NMSE fitness for one candidate at one site.
\end{itemize}

\vspace{0.2cm}
\textbf{Why this split?}
\begin{itemize}
  \item The search logic is generic.
  \item Whisper and Parakeet differ mainly in where rotations are inserted and how a block is replayed.
  \item This makes it easy to add more models later without rewriting the search loop.
\end{itemize}
\end{frame}

% 21
\begin{frame}{Methodology: common site abstraction}
\textbf{First shared site type:} attention value/output rotation.

\vspace{0.2cm}
For both Whisper and Parakeet, define a site as:
\begin{itemize}
  \item a block index $\ell$,
  \item the value/output rotation parameters for that block,
  \item the candidate family
  \[
  R_2^{(\ell)} = D_0^{(\ell)} H D_1^{(\ell)}.
  \]
\end{itemize}

\vspace{0.2cm}
\textbf{Common site API idea:}
\begin{itemize}
  \item initialize a candidate population,
  \item mutate a candidate,
  \item temporarily apply a candidate to the block,
  \item score the resulting fake-quantized block output with NMSE.
\end{itemize}
\end{frame}

% 22
\begin{frame}{Methodology: caching for efficient local search}
\textbf{Cache once, search many times.}

\vspace{0.2cm}
For each searchable block, cache:
\begin{itemize}
  \item the block input on a small calibration subset,
  \item all masks and auxiliary tensors needed to replay the block exactly,
  \item the floating-point block output
  \[
  Y_{\ell}^{\mathrm{FP}}.
  \]
\end{itemize}

\vspace{0.2cm}
Then candidate evaluation only needs:
\begin{itemize}
  \item insert the candidate rotation at that site,
  \item replay the block on cached inputs under fake quantization,
  \item compute
  \[
  \frac{\|Y_{\ell}^{\mathrm{FP}}-Y_{\ell}^{\mathrm{fake\text{-}quant}}\|_2^2}
  {\|Y_{\ell}^{\mathrm{FP}}\|_2^2+\varepsilon}.
  \]
\end{itemize}

\vspace{0.2cm}
This avoids re-running the full ASR model for every offspring candidate.

\vspace{0.2cm}
\textbf{First implementation choice:}
\begin{itemize}
  \item Cache only block inputs, required auxiliary tensors, and floating-point block outputs.
  \item Do \emph{not} cache deeper internal attention-path tensors in the first version.
  \item This keeps the replay logic simple and ensures the local score reflects the true block behavior.
  \item In practice, the adapter should store the full block-call context, i.e. the input tensor plus replay kwargs, rather than only the main hidden-state tensor.
\end{itemize}
\end{frame}

% 12
% \begin{frame}{Positioning and expected contribution}
% \begin{center}
% \begin{tabular}{>{\raggedright\arraybackslash}p{0.26\linewidth} >{\raggedright\arraybackslash}p{0.30\linewidth} >{\raggedright\arraybackslash}p{0.26\linewidth}}
% \toprule
% Method & Rotation family & Optimization cost \\
% \midrule
% QuaRot-style & Random Hadamard & Very low \\
% SpinQuant & Dense orthonormal matrices & High memory / gradient-based \\
% Proposed direction & Searched signed-Hadamard family & Low-memory / gradient-free \\
% \bottomrule
% \end{tabular}
% \end{center}

% \vspace{0.3cm}
% The goal is not necessarily to beat SpinQuant in absolute accuracy.

% \vspace{0.2cm}
% A strong contribution would be:
% \[
% \text{recover most of SpinQuant's gain with QuaRot-like efficiency.}
% \]

% \vspace{0.2cm}
% Important metrics:
% \[
% \text{accuracy gap},\quad \text{search time},\quad \text{peak GPU memory},\quad \text{inference overhead}.
% \]
% \end{frame}

% 13
% \begin{frame}{Takeaway}
% \vspace{0.2cm}
% \textbf{SpinQuant:}
% \[
% \text{Learn rotations to make tensors easier to quantize.}
% \]

% \vspace{0.3cm}
% \textbf{EvoPress:}
% \[
% \text{Use efficient search when compression choices interact globally.}
% \]

% \vspace{0.3cm}
% \textbf{Proposed idea:}
% \[
% \text{Search over fast Hadamard-like rotations for outlier reduction.}
% \]

% \vspace{0.4cm}
% This gives a clean research direction:
% \[
% \boxed{\text{Adaptive rotations without dense rotation learning.}}
% \]
% \end{frame}

% 14
% \begin{frame}{References}
% \small
% \begin{itemize}
%   \item Liu et al., \textit{SpinQuant: LLM Quantization with Learned Rotations}, ICLR 2025.
%   \item Sieberling et al., \textit{EvoPress: Accurate Dynamic Model Compression via Evolutionary Search}, ICML 2025.
%   \item Ashkboos et al., \textit{QuaRot: Outlier-Free 4-Bit Inference in Rotated LLMs}.
% \end{itemize}
% \end{frame}

\end{document}
